\newcounter{entries}
\newcommand{\educationtext}[3]{%
{\normalsize\textbf{#1}}\\[0.1em]
\textit{#2}\\[0.05em]
{\fontsize{9.5}{12}\selectfont
\setlength{\parskip}{0.4em}
#3
}%
}
\newcommand{\education}[4]{%
    \stepcounter{entries}
    \pgfmathtruncatemacro{\itemno}{\value{entries}}
    \node[
        below=\sectionsep of #4.south west,
        eventdottext]
    (item \itemno header)
    {\phantom{#1}};
    \node[
        below=\sectionsep of #4.south west,
        eventdottext
    ](item \itemno){\educationtext{#1}{#2}{#3}};
    \node[
        left=\timedotsep of item \itemno header,
        timedot
    ](time \itemno)
    {};
}
\newcommand{\experiencetext}[4]{%
{\normalsize\textbf{#1, #2}}\\[0.1em]
\textit{#3}\\[0.05em]
{\fontsize{9.5}{12}\selectfont
\setlength{\parskip}{0.4em}
#4
}%
}
\newcommand{\experience}[5]{%
    \stepcounter{entries}
    \pgfmathtruncatemacro{\itemno}{\value{entries}}
    \node[
        below=\sectionsep of #5.south west,
        eventdottext]
    (item \itemno header)
    {\phantom{#1}};
    \node[
        below=\sectionsep of #5.south west,
        eventdottext
    ](item \itemno){\experiencetext{#1}{#2}{#3}{#4}};
    \node[
        left=\timedotsep of item \itemno header,
        timedot
    ](time \itemno)
    {};
}
\newcommand{\sectionstroke}[1]{%
    \begin{scope}[on background layer]
        \draw[line width=1.5pt,sidebar]
        let \p1=(#1.south west),
        \p2=(current page.east) in
        (\x1,\y1-\sectionheaderbelow) to (\x2,\y1-\sectionheaderbelow);
    \end{scope}
}

\setlength{\sectionsep}{1.2em}

\begin{tikzpicture}[
        every node/.style={
                inner sep=0pt,
                outer sep=0pt
            },
        remember picture,
        overlay,
        shift={($(current page.north west)%
                    +(\sidewidth+2.75\margin+\timedotsep,-0.95\margin)$)}]
    \hypersetup{urlcolor=linkonlight}

    \node[sectionTitle] at (0,-0) (title 1) {\cvsection{Experience}};
    \node[left=\timedotsep of title 1,headerIcon] {\faBriefcase};
    \sectionstroke{title 1}
    \experience{Developer}{Acturis Ltd}{%
        August 2024 - present
    }{%
        Develop predominantly in C\# on the dual-entity quoting and documents
        engine underpinning the Acturis application.

        Lead the design, implementation and rollout of large projects, working
        with client-facing business analysts and other developers to gather
        requirements, define technical solutions, manage scope and coordinate
        delivery from initial specification through to implementation.

        Own technical work across the full development lifecycle, contributing
        to architectural and design decisions while balancing business
        requirements, technical constraints and delivery timelines.

        Mentor and support other developers through onboarding, technical guidance,
        code reviews and knowledge sharing, helping colleagues develop their technical
        skills and confidence to take on increasingly complex work.

        Contribute to the improvement of CI/CD processes, increasing automation
        across the development workflow and reducing manual engineering effort.
    }{title 1}
    \experience{Research Intern}{Huawei R\&D, Edinburgh}{%
        January 2021 - November 2023
    }{%
        Worked on the compiler for Huawei's
        open-source Cangjie programming language, written in C++ against an LLVM
        backend.

        Led the development of a domain-specific language for hardware design
        and reasoning, taking the project from initial research and requirements
        through to language design, implementation, and evaluation.

        Developed and implemented compiler optimisations for the new language,
        solving complex problems in compiler engineering, language
        implementation, and code generation.
    }{item 1}
    \experience{Teaching Assistant}{University of Birmingham}{%
        September 2019 - June 2024
    }{%
        Led teaching across the \textit{Compilers \& Languages},
        \textit{Logical and Mathematical Foundations for Computer Science} and
        \textit{Theories of Computation} modules, delivering tutorials for groups
        of 10-20 students.

        Coordinated assignment marking across teaching assistants, supporting
        consistent assessment and helping develop teaching materials and
        approaches.

        Received awards in 2022 for going above and beyond my teaching
        responsibilities.
    }{item 2}
    \node[
        left=\timedotsep of item 3.south west,
        invisibletimedot
    ](time 3)
    {};
    \draw (time 1.center) -- (time 3.center);

    %%%%%%%%%%%
    % Education
    %%%%%%%%%%%

    \node[below=0.65cm of item 3.south west,sectionTitle] (title 2) {\cvsection{Education}};
    \node[left=\timedotsep of title 2,headerIcon] {\faGraduationCap};
    \node[below=0.45cm of item 3.south west,sectionTitle] (title 1 dummy) {\phantom{\cvsection{Education}}};
    \sectionstroke{title 2}
    \education{%
        PhD in Computer Science, University of Birmingham%
    }{%
        September 2019 - July 2024%
    }{%
        Research in foundations of digital circuits, applying techniques from
        programming languages and compiler research to hardware optimisation and
        verification.

        Designed and implemented research prototypes in OCaml and ReScript for
        experimentation and evaluation, and presented technical work to
        academic and industrial audiences at international conferences.
    }{title 2}%
    \education{%
        MSci in Computer Science, University of Birmingham (Honours, Class I)%
    }{%
        September 2015 - July 2019
    }{%
        Relevant modules included Compilers \& Languages,
        Functional Programming, Computer-Aided Verification, C/C++, and
        Operating Systems.
    }{item 4}%

    \node[
        left=\timedotsep of item 5.south west,
        invisibletimedot
    ](time 5)
    {};
    \draw (time 4.center) to (time 5.center);

    %%%%%%
    % Volunteering
    %%%%%%%

    \node[below=0.65cm of item 5.south west,sectionTitle] (title 3) {\cvsection{Volunteering}};
    \node[left=\timedotsep of title 3,headerIcon] {\faRocket};
    \node[below=0.45cm of item 3.south west,sectionTitle] (title 3 dummy) {\phantom{\cvsection{Volunteering}}};
    \sectionstroke{title 3}
    \education{Local Organiser, Twelfth Symposium on Compositional Structures}{%
        15-16 April 2024}{
        Helped organise and host the conference in Birmingham, securing funding
        and coordinating the website, social programme, catering, venues, event
        setup and talk selection with the organising committee.
    }{title 3}
    \node[
        left=\timedotsep of item 6.south west,
        invisibletimedot
    ](time 9){};
    \draw (time 6.center) -- (time 9.center);

\end{tikzpicture}
